%---------------------------------------------------------------------
\chapter{Agentic Concepts Demonstrated}
%---------------------------------------------------------------------
This chapter maps the system to the concept checklist of the course guidelines. Each tick points
to code, a test or a recorded run; Appendix~\ref{app:traces} reproduces trace excerpts.

\section{Concept Coverage Matrix}
\begin{table}[htbp]
\centering
\caption{Coverage of the agentic-concept checklist}
\label{tab:concepts}
\small
\begin{tabularx}{\linewidth}{@{}p{1.8cm} >{\raggedright\arraybackslash}p{3.7cm} >{\centering\arraybackslash}p{1.2cm} >{\raggedright\arraybackslash}X@{}}
\toprule
\textbf{Tier} & \textbf{Concept} & \textbf{Status} & \textbf{Evidence} \\ \midrule
Mandatory & Tool / function use & \checkmark & Six audited tools (Table~\ref{tab:tools}); \texttt{tests/\allowbreak integration/\allowbreak test\_\allowbreak a07\_\allowbreak tools.py} \\
Mandatory & Agentic RAG & \checkmark & Graded \texttt{search\_rules}; MG2 Retriever with re-query; \texttt{tests/\allowbreak unit/\allowbreak test\_\allowbreak h06\_\allowbreak retrieval.py}, \texttt{test\_\allowbreak agents.py} \\
Mandatory & Planning and reasoning & \checkmark & \textsc{plan}/\textsc{assess} loop, Section~4.4; \texttt{tests/\allowbreak integration/\allowbreak test\_\allowbreak a06\_\allowbreak workflow.py} \\
Mandatory & Observability / tracing & \checkmark & Ledger, result records and audit chain, Section~7.1 \\
Mandatory & Evaluation & \checkmark & Chapter~7; 1{,}008 scored trials in the MG2 series alone, plus MG1 \\
Bonus & Multi-agent collaboration & \checkmark & Retriever and Classifier (MG2), measured; result negative \\
Bonus & Memory & Partial & Run state checkpointed; no memory across cases (by design) \\
Bonus & Human-in-the-loop and guardrails & \checkmark & Finding gate, review bound to revision; \texttt{tests/\allowbreak unit/\allowbreak test\_\allowbreak finding\_\allowbreak gate.py}, \texttt{test\_\allowbreak v05\_\allowbreak review.py} \\
Bonus & Durable workflow & \checkmark & SQLite checkpoints, CAS versions, resume; \texttt{tests/\allowbreak integration/\allowbreak test\_\allowbreak cp4\_\allowbreak recovery.py} \\
Bonus & Experimentation & \checkmark & Pre-registered MG2 experiment and held-out run, Chapter~8 \\
\bottomrule
\end{tabularx}
\end{table}

\section{Evidence for Each Concept}

\subsection{Tool Use}
The model never touches data directly. In \textsc{plan} it names one tool from the allowed list with
typed arguments; the runner validates the arguments against the case (an observation identifier
that does not exist is refused and the model is asked once more), records the intent, executes the
tool and returns its structured result. In the trial records, every result carries a
\texttt{calls} list in which model calls and stage actions alternate, and a budget block counting
actions, tool calls and model calls (Listing~\ref{lst:result} in Appendix~\ref{app:traces}).

\subsection{Agentic RAG}
\texttt{search\_rules} queries the BM25 index over the two regulatory documents, filters
provisions by temporal validity and grades what it returns before the agent sees it. In MG2 the
Retriever applies the corrective pattern explicitly: after each read or search it grades the result
as \texttt{SUFFICIENT}, \texttt{WEAK} or \texttt{EMPTY}, and after a weak or empty result it must
change its request, because an identical request is refused. Its claims are then checked by a
deterministic quote verifier. Chapter~8 reports how this behaved on the real model.

\subsection{Planning and Reasoning}
Each turn separates the decision to act from the judgement of what the action showed. The
\textsc{assess} step must return one of four verdicts and three dispositions, and \texttt{RESOLVE}
is allowed only with \texttt{SUPPORTS}. The adaptive variant (V2) chooses its own checks; the fixed
variant (V1) runs a predeclared sequence and only assesses, which allows the value of the planning
step to be measured directly.

\subsection{Observability}
Every trial writes a ledger line when it starts and when it ends (status, wall time, tokens and the
digests of its bundle and result), and a result record with the run status, stop reason, call list,
budget, evidence census, findings, rejected proposals and the verdict of the audit-chain check.
Section~7.1 describes what is captured and how it is used.

\subsection{Evaluation}
Trials are scored by a deterministic evaluator against hidden answers, with model-free controls in
every session and case-level bootstrap intervals (Chapter~7).

\subsection{Bonus Concepts}
\textbf{Multi-agent collaboration.} The MG2 Retriever and Classifier were built, tested with
scripted model ports, and then measured on the real model with three prompt versions in the T4 series
(1{,}008 trials including baselines and controls); Chapter~8 reports the negative result.
\textbf{Human in the loop and guardrails.} Model proposals pass a deterministic gate; the reviewer's
decision binds the exact report revision; dossier text is treated as untrusted data, and one
fault family (F16) plants instructions in the dossier to test this.
\textbf{Durable workflow.} Graph state is checkpointed with compare-and-swap versions, and an
interrupted run resumes without repeating a dispatched action; the evaluation queue also resumes
every session from its own ledger.
\textbf{Memory} is limited to the run: by design, nothing learned on one case is carried to another,
because that would leak information between evaluation cases.
